\section*{Exercice Simplexe : Routage des camions d'eau}

\subsection*{Contexte}
Une société organise la distribution d’eau vers deux zones périphériques.  
Chaque tournée permet de livrer une quantité d’eau différente selon la zone.

\subsection*{Variables de décision}
\[
X = \text{nombre de tournées vers la zone 1}
\]

\[
Y = \text{nombre de tournées vers la zone 2}
\]

\subsection*{Fonction objectif}
L’objectif est de maximiser la quantité totale d’eau distribuée :

\[
\max Z = 30X + 40Y
\]

\subsection*{Contraintes}
\[
X + Y \leq 50
\]

\[
3X + 4Y \leq 120
\]

\[
2X + 5Y \leq 150
\]

\[
X \geq 0, \quad Y \geq 0
\]

\subsection*{Forme standard}
On ajoute les variables d'écart \(S_1\), \(S_2\) et \(S_3\) :

\[
X + Y + S_1 = 50
\]

\[
3X + 4Y + S_2 = 120
\]

\[
2X + 5Y + S_3 = 150
\]

\[
Z - 30X - 40Y = 0
\]

\subsection*{Résolution}
Les sommets réalisables principaux sont :

\[
(0,0), \quad (40,0), \quad (0,30)
\]

Calcul de la fonction objectif :

\[
Z(0,0) = 0
\]

\[
Z(40,0) = 30(40) + 40(0) = 1200
\]

\[
Z(0,30) = 30(0) + 40(30) = 1200
\]

\subsection*{Solution optimale}
On obtient plusieurs solutions optimales.  
Par exemple :

\[
X = 40, \quad Y = 0
\]

ou :

\[
X = 0, \quad Y = 30
\]

Dans les deux cas :

\[
Z_{\max} = 1200
\]

\subsection*{Interprétation}
La société peut choisir plusieurs plans de tournées donnant la même quantité maximale distribuée.  
Cela signifie qu’il existe plusieurs organisations possibles avec le même résultat optimal.